% GENERATED FILE. Edit canonical lessons, then run npm run generate:books.
% canonical-source-hash: fnv1a64:0978e6080df10841
% canonical-lessons: PT-C315-funeral, PT-C315-enterrar, PT-C315-enterro, PT-C315-cemiterio, PT-C315-tumulo

\chapter{O funeral}
\label{ch:pt-o-funeral}
\hlchaptermodality{315}

\begin{chapteropening}
\textbf{By the end of this chapter:} \emph{I can talk about a funeral, a burial and the cemetery.}

Complete the last lesson of chapter 315: I can talk about a funeral, a burial and the cemetery.
\end{chapteropening}

\section[o funeral]{o funeral — the funeral}

\label{lesson:PT-C315-funeral}



\begin{quote}
Before the new one: say the Portuguese for mourning, then the Portuguese for the wake.
\end{quote}

\subsection*{You'll want to know: o funeral}
\textbf{o funeral} — \textquotedblleft{}the funeral\textquotedblright{}.

\begin{quote}
O funeral foi na sexta-feira de manhã.
\end{quote}

\begin{quote}
\emph{The funeral was on Friday morning.}
\end{quote}

\emph{Ir ao funeral} — to go to the funeral.

\begin{cousinweb}
\emph{funeral} was borrowed from Late Latin \emph{fūnerālis} \textquotedblleft{}of a funeral\textquotedblright{} (EtymDB), from \emph{fūnus} \textquotedblleft{}a funeral, death\textquotedblright{}.
\end{cousinweb}

\begin{grammarlens}[title={What you've built}]
One more word for telling, explaining and describing at B1.
\end{grammarlens}

\subsection*{Your turn}
\begin{itemize}
\raggedright
  \item \emph{Say it:} \emph{o funeral}
  \item \emph{Say it:} \emph{o funeral}, once more
  \item \emph{Say it:} a sentence of your own with \emph{o funeral}
  \item \emph{Recall:} say the Portuguese for mourning, then the Portuguese for the wake, then say \emph{o funeral} again
\end{itemize}

\begin{tcolorbox}[breakable,colback=teal!4,colframe=teal!35!black,title={Before you move on}]
What does o funeral mean? (\textquotedblleft{}The funeral\textquotedblright{}.) Say it once more.
\end{tcolorbox}
\section[enterrar]{enterrar — to bury}

\label{lesson:PT-C315-enterrar}



\begin{quote}
Before the new one: say the Portuguese for the wake, then the Portuguese for the funeral.
\end{quote}

\subsection*{You'll want to know: enterrar}
\textbf{enterrar} — \textquotedblleft{}to bury\textquotedblright{}.

\begin{quote}
O cão enterrou o osso no quintal.
\end{quote}

\begin{quote}
\emph{The dog buried the bone in the yard.}
\end{quote}

For people and for things.

\begin{cousinweb}
\emph{enterrar} is inherited from Vulgar Latin \emph{interrāre} \textquotedblleft{}to put in the earth\textquotedblright{} (EtymDB), built on \emph{terra} \textquotedblleft{}earth\textquotedblright{}.
\end{cousinweb}

\begin{grammarlens}[title={What you've built}]
One more word for telling, explaining and describing at B1.
\end{grammarlens}

\subsection*{Your turn}
\begin{itemize}
\raggedright
  \item \emph{Say it:} \emph{enterrar}
  \item \emph{Say it:} \emph{enterrar}, once more
  \item \emph{Say it:} a sentence of your own with \emph{enterrar}
  \item \emph{Recall:} say the Portuguese for the wake, then the Portuguese for the funeral, then say \emph{enterrar} again
\end{itemize}

\begin{tcolorbox}[breakable,colback=teal!4,colframe=teal!35!black,title={Before you move on}]
What does enterrar mean? (\textquotedblleft{}To bury\textquotedblright{}.) Say it once more.
\end{tcolorbox}
\section[o enterro]{o enterro — the burial}

\label{lesson:PT-C315-enterro}



\begin{quote}
Before the new one: say the Portuguese for the funeral, then the Portuguese for to bury.
\end{quote}

\subsection*{You'll want to know: o enterro}
\textbf{o enterro} — \textquotedblleft{}the burial\textquotedblright{}.

\begin{quote}
Depois da missa, houve o enterro.
\end{quote}

\begin{quote}
\emph{After the mass came the burial.}
\end{quote}

Built on \emph{enterrar}.

\begin{cousinweb}
\emph{enterro} is built on \emph{enterrar} \textquotedblleft{}to bury\textquotedblright{} (EtymDB), from Vulgar Latin \emph{interrāre} \textquotedblleft{}to put in the earth\textquotedblright{}, from \emph{terra}.
\end{cousinweb}

\begin{grammarlens}[title={What you've built}]
One more word for telling, explaining and describing at B1.
\end{grammarlens}

\subsection*{Your turn}
\begin{itemize}
\raggedright
  \item \emph{Say it:} \emph{o enterro}
  \item \emph{Say it:} \emph{o enterro}, once more
  \item \emph{Say it:} a sentence of your own with \emph{o enterro}
  \item \emph{Recall:} say the Portuguese for the funeral, then the Portuguese for to bury, then say \emph{o enterro} again
\end{itemize}

\begin{tcolorbox}[breakable,colback=teal!4,colframe=teal!35!black,title={Before you move on}]
What does o enterro mean? (\textquotedblleft{}The burial\textquotedblright{}.) Say it once more.
\end{tcolorbox}
\section[o cemitério]{o cemitério — the cemetery}

\label{lesson:PT-C315-cemiterio}



\begin{quote}
Before the new one: say the Portuguese for to bury, then the Portuguese for the burial.
\end{quote}

\subsection*{You'll want to know: o cemitério}
\textbf{o cemitério} — \textquotedblleft{}the cemetery\textquotedblright{}.

\begin{quote}
O cemitério fica ao lado da igreja.
\end{quote}

\begin{quote}
\emph{The cemetery is next to the church.}
\end{quote}

Every village has its own.

\begin{cousinweb}
\emph{cemitério} was borrowed from Late Latin \emph{coemētērium} (EtymDB), from Greek \emph{koimēt\'{\={e}}rion} \textquotedblleft{}a sleeping place\textquotedblright{}.
\end{cousinweb}

\begin{grammarlens}[title={What you've built}]
One more word for telling, explaining and describing at B1.
\end{grammarlens}

\subsection*{Your turn}
\begin{itemize}
\raggedright
  \item \emph{Say it:} \emph{o cemitério}
  \item \emph{Say it:} \emph{o cemitério}, once more
  \item \emph{Say it:} a sentence of your own with \emph{o cemitério}
  \item \emph{Recall:} say the Portuguese for to bury, then the Portuguese for the burial, then say \emph{o cemitério} again
\end{itemize}

\begin{tcolorbox}[breakable,colback=teal!4,colframe=teal!35!black,title={Before you move on}]
What does o cemitério mean? (\textquotedblleft{}The cemetery\textquotedblright{}.) Say it once more.
\end{tcolorbox}
\section[o túmulo]{o túmulo — the tomb, the grave}

\label{lesson:PT-C315-tumulo}



\begin{quote}
Before the new one: say the Portuguese for the burial, then the Portuguese for the cemetery.
\end{quote}

\subsection*{You'll want to know: o túmulo}
\textbf{o túmulo} — \textquotedblleft{}the tomb, the grave\textquotedblright{}.

\begin{quote}
No túmulo havia flores frescas.
\end{quote}

\begin{quote}
\emph{There were fresh flowers on the grave.}
\end{quote}

\emph{A campa} is a simpler grave in the ground.

\begin{cousinweb}
\emph{túmulo} is a learned borrowing from Latin \emph{tumulus} \textquotedblleft{}a mound\textquotedblright{} (EtymDB), from \emph{tumēre} \textquotedblleft{}to swell\textquotedblright{}.
\end{cousinweb}

\begin{grammarlens}[title={What you've built}]
One more word for telling, explaining and describing at B1.
\end{grammarlens}

\subsection*{Your turn}
\begin{itemize}
\raggedright
  \item \emph{Say it:} \emph{o túmulo}
  \item \emph{Say it:} \emph{o túmulo}, once more
  \item \emph{Say it:} a sentence of your own with \emph{o túmulo}
  \item \emph{Recall:} say the Portuguese for the burial, then the Portuguese for the cemetery, then say \emph{o túmulo} again
\end{itemize}

\begin{tcolorbox}[breakable,colback=teal!4,colframe=teal!35!black,title={Before you move on}]
What does o túmulo mean? (\textquotedblleft{}The tomb, the grave\textquotedblright{}.) Say it once more.
\end{tcolorbox}
